% 原创教学示意；分三步显示前向依赖、单路径回传和多路径累加。
\begin{tikzpicture}[x={\dimexpr\linewidth/169\relax},y=1mm,
  font=\figurefont\footnotesize,text=NNInk,
  small op/.style={nn operator,minimum width=13mm,minimum height=10mm},
  state/.style={nn operator,minimum width=14mm,minimum height=10mm}]
  \path[fill=FigureBlueFill,rounded corners=1mm] (0,-27) rectangle (53,32);
  \path[fill=FigureYellowFill,rounded corners=1mm] (58,-27) rectangle (111,32);
  \path[fill=FigureRedFill,rounded corners=1mm] (116,-27) rectangle (169,32);
  \node[nn heading] at (26.5,26) {(a) 前向：一份状态被两处读取};
  \node[nn heading] at (84.5,26) {(b) 单条路径回传};
  \node[nn heading] at (142.5,26) {(c) 多路径贡献相加};

  \begin{scope}[shift={(1,0)}]
    \node[state] (a-sender) at (8,0) {$\bm h_j^{(l)}$};
    \foreach \idx/\yy in {i/11,k/-11}{
      \node[small op] (a-m-\idx) at (29,\yy) {$M_{\,\idx\leftarrow j}$};
      \node[state] (a-state-\idx) at (47,\yy) {$\bm h_{\idx}^{(l+1)}$};
      \draw[nn flow,draw=NNHidden] (a-sender.east)--(a-m-\idx.west);
      \draw[nn flow,draw=NNHidden] (a-m-\idx.east)--(a-state-\idx.west);
    }
  \end{scope}

  \begin{scope}[shift={(59,0)}]
    \node[state] (b-sender) at (8,0) {$\bm h_j^{(l)}$};
    \node[small op] (b-message) at (29,0) {$M_{\,i\leftarrow j}$};
    \node[state] (b-state) at (47,0) {$\bm h_i^{(l+1)}$};
    \draw[nn flow,draw=NNHidden] (b-sender.east)--(b-message.west);
    \draw[nn flow,draw=NNHidden] (b-message.east)--(b-state.west);
    \draw[nn gradient,line width=1.3pt]
      ([yshift=-2.2mm]b-state.west)--([yshift=-2.2mm]b-message.east);
    \draw[nn gradient,line width=1.3pt]
      ([yshift=-2.2mm]b-message.west)--([yshift=-2.2mm]b-sender.east);
    \node[text=FigureInk] at (29,-14) {$\bm\delta_{i\leftarrow j}^{(l)}$};
  \end{scope}

  \begin{scope}[shift={(116,0)}]
    \node[state] (c-sender) at (7,0) {$\bm h_j^{(l)}$};
    \node[nn pointwise,draw=NNGradient] (c-sum) at (20,0) {$+$};
    \foreach \idx/\yy in {i/11,k/-11}{
      \node[small op] (c-m-\idx) at (29,\yy) {$M_{\,\idx\leftarrow j}$};
      \node[state] (c-state-\idx) at (42,\yy) {$\bm h_{\idx}^{(l+1)}$};
      \draw[nn flow,draw=NNHidden,sharp corners]
        (c-sender.east)--(14,0)--(14,\yy)--(c-m-\idx.west);
      \draw[nn flow,draw=NNHidden] (c-m-\idx.east)--(c-state-\idx.west);
      \draw[nn gradient,line width=1.3pt]
        ([yshift=-2.2mm]c-state-\idx.west)--
        ([yshift=-2.2mm]c-m-\idx.east);
      \draw[nn gradient,line width=1.3pt]
        ([yshift=-2.2mm]c-m-\idx.west)--(c-sum.east);
    }
    \draw[nn gradient,line width=1.5pt] (c-sum.west)--(c-sender.east);
    \node[text=FigureInk,align=center] at (26,-22)
      {$\bm g_j^{(l)}=\bm\delta_{i\leftarrow j}^{(l)}
      +\bm\delta_{k\leftarrow j}^{(l)}$};
  \end{scope}
\end{tikzpicture}
